\documentclass[11pt,a4paper]{ctexart}

\usepackage[
  left=1.55cm,
  right=1.55cm,
  top=0.75cm,
  bottom=0.85cm
]{geometry}
\usepackage{tabularx}
\usepackage{enumitem}
\usepackage{titlesec}
\usepackage{xcolor}
\usepackage{graphicx}
\usepackage{fontawesome5}
\usepackage[colorlinks=true,linkcolor=blue,urlcolor=blue]{hyperref}
\usepackage{array}

\pagestyle{empty}
\setlength{\parindent}{0pt}
\setlength{\parskip}{0pt}
\emergencystretch=1.5em

\definecolor{main}{RGB}{35,35,35}
\definecolor{sub}{RGB}{95,95,95}

% ==================== 字体层级 ====================
% 尽量避免过重的 \bfseries，主要用字号、黑体和灰度区分层级
\newcommand{\namefont}{\fontsize{29pt}{32pt}\selectfont\heiti}
\newcommand{\contactfont}{\fontsize{11.8pt}{13.4pt}\selectfont\songti\color{main}}
\newcommand{\sectiontitlefont}{\fontsize{12.8pt}{14.2pt}\selectfont\heiti\color{main}}
\newcommand{\entrytitlefont}{\fontsize{10.8pt}{12.2pt}\selectfont\heiti\color{main}}
\newcommand{\bodyfont}{\fontsize{9.8pt}{11.8pt}\selectfont\songti\color{main}}
\newcommand{\metafont}[1]{{\fontsize{9.5pt}{11pt}\selectfont\songti\color{sub}#1}}
\newcommand{\tagfont}[1]{{\fontsize{9.5pt}{11pt}\selectfont\songti\color{sub}#1}}
\newcommand{\awardfont}{\fontsize{9.8pt}{11.6pt}\selectfont\songti\color{main}}

% ==================== Section 样式 ====================
\titleformat{\section}
  {\sectiontitlefont}
  {}
  {0pt}
  {}
  [\vspace{-0.42em}{\color{main}\rule{\linewidth}{0.45pt}}\vspace{-0.75em}]

\titlespacing*{\section}{0pt}{0.55em}{0.35em}

% ==================== 列表样式 ====================
\setlist[itemize,1]{
  label={\raisebox{0.2ex}{\scriptsize$\bullet$}},
  leftmargin=1.3em,
  itemsep=0.08em,
  topsep=0.08em,
  parsep=0em,
  before=\bodyfont
}

\setlist[itemize,2]{
  label=\textendash,
  leftmargin=2.35em,
  itemsep=0.04em,
  topsep=0.02em,
  parsep=0em,
  before=\bodyfont
}

% ==================== 基础命令 ====================
\newcommand{\eduentry}[4]{
  \noindent
  {\entrytitlefont #1} \quad {\bodyfont #2} \quad {\bodyfont #3}
  \hfill \metafont{#4}
  \par\vspace{0em}
}

\newcommand{\edulabel}[1]{{\heiti #1：}}

% 右上角照片：请确保 sjx.jpg 与 tex 文件在同一目录下
\newcommand{\photobox}{
  {\includegraphics[width=3.6cm,height=4.8cm,keepaspectratio]{sjx.jpg}}
}

\newcommand{\resumeheader}[3]{
  \noindent
  \begin{tabular}{@{}m{3.8cm}@{}>{\centering\arraybackslash}m{\dimexpr\linewidth-7.6cm\relax}@{}>{\raggedleft\arraybackslash}m{3.8cm}@{}}
    &
    \begin{minipage}{\linewidth}
      \centering
      {\namefont #1}\\[0.15em]
      {\contactfont #2 \quad | \quad \href{mailto:#3}{\textcolor{main}{\underline{#3}}}}
    \end{minipage}
    &
    \hspace{-0.45cm}\photobox
  \end{tabular}
  \vspace{0em}
}

\newcommand{\projecttitle}[1]{{\entrytitlefont #1}}

\newcommand{\project}[3]{
  \noindent
  \projecttitle{#1} {\bodyfont\color{main}{\color{sub}\,|\,}#3}
  \hfill \metafont{#2}\par
  \vspace{-0.02em}
}

\newcommand{\projectwithlink}[4]{
  \noindent
  \begin{tabular*}{\linewidth}{@{}l@{\extracolsep{\fill}}r@{}}
    \projecttitle{#1} {\bodyfont\color{main}{\color{sub}\,|\,}#3{\color{sub}\,|\,}\href{#4}{\textcolor{main}{\raisebox{-0.08em}{\faGithub}\,GitHub}}} &
    \metafont{#2} \\
  \end{tabular*}\par
  \vspace{-0.02em}
}

% 与 projectwithlink 同一排版；第 4、5 个参数分别为 GitHub 与 arXiv URL。
% arXiv URL 为空时显示不可点击的 arXiv 标签。
\newcommand{\projectwithpaper}[5]{
  \noindent
  \begin{tabular*}{\linewidth}{@{}l@{\extracolsep{\fill}}r@{}}
    \projecttitle{#1} {\bodyfont\color{main}{\color{sub}\,|\,}#3%
      {\color{sub}\,|\,}\href{#4}{\textcolor{main}{\raisebox{-0.08em}{\faGithub}\,GitHub}}%
      {\color{sub}\,|\,}%
      \if\relax\detokenize{#5}\relax
        \textcolor{sub}{\raisebox{-0.08em}{\faFilePdf}\,arXiv}%
      \else
        \href{#5}{\textcolor{main}{\raisebox{-0.08em}{\faFilePdf}\,arXiv}}%
      \fi
    } &
    \metafont{#2} \\
  \end{tabular*}\par
  \vspace{-0.02em}
}

\newcommand{\compactproject}[2]{
  \noindent
  \projecttitle{#1} \hfill \metafont{#2}\par
  \vspace{-0.02em}
}

% 仿照泽子 CV 的获奖部分：无 bullet，紧凑逐行，年份靠右
\newcommand{\awarditem}[2]{
  \noindent
  {\awardfont #1 \hfill {\color{sub}#2}}\par
  \vspace{0.02em}
}

\begin{document}

\resumeheader
{孙家兴}
{+86 188-1171-7335}
{2501213407@stu.pku.edu.cn}

\vspace{0.12em}

\section{教育背景}

\eduentry
{北京大学}
{智能学院}
{硕士研究生（三年学硕）}
{2025.09 -- 至今}

\vspace{0.10em}

\eduentry
{北京大学}
{信息科学技术学院}
{本科}
{2021.09 -- 2025.06}

\section{研究经历}

\projectwithlink
{熵自适应策略蒸馏}
{2026.07 -- 至今}
{Python / PyTorch / veRL / ALFWorld}
{https://github.com/Sisyphe-lee/opd-baseline-repro}

\begin{itemize}
  \item 提出熵自适应蒸馏，依据教师熵动态选择可靠轨迹前缀。
  \item ALFWorld full274 成功率 86.86\%，较 TCOD-F2B / Vanilla OPD 提升 2.19 / 7.30 个百分点。
\end{itemize}

\vspace{0.2em}

% TODO：论文发布后，在最后一个参数中填入 arXiv URL。
\projectwithpaper
{Power Sampling 失效分析与修复}
{2026.06 -- 2026.08}
{Power Sampling / Self-Consistency / SoftSat}
{https://github.com/jiaxingsunpku/SoftSat}
{}

\begin{itemize}
  \item 揭示 Power Sampling 可增加正确轨迹质量却降低多轨迹决策性能，并将失效归因于 coverage mismatch 与 dose mismatch；提出 Relative-Rank SoftSat，以相对秩校准锐化强度、以饱和增益保留中等概率路径。
  \item 在 3 个模型、3 个代码/数学/物理基准的 9 组设置中，固定指数 Power 在 7 组退化、最大下降 18.5 个百分点；SoftSat 在 5 组提升、1 组持平，最大回退仅 1.0 个百分点，加权实现达到 99.55\% 决策一致与 7.42$\times$ 推理加速。
\end{itemize}

\section{项目经历}

\projectwithlink
{LightCC Coder: Code Agent Harness}
{2026.05 -- 至今}
{TypeScript / Agent Loop / Repo Tooling}
{https://github.com/Sisyphe-lee/light-cc-coder/tree/main}

\begin{itemize}
  \item 参与构建面向真实代码仓库的轻量 Claude Code 风格 coder harness，核心约 10k 行 TypeScript；实现交互式 REPL、one-shot CLI、OpenAI-compatible provider、JSONL 事件流，以及受限仓库工具链、transcript 回放、上下文压缩、resume 与可选 OS sandbox。
  \item 使用 deepseek-v4-flash 在固定子集评测：SWE-bench Verified resolved 10/20、Terminal-Bench 2.1 通过 11/20；SWE-bench 横评中平均请求数、成本、token 使用量和运行时间均低于 OpenCode / OpenHands。
\end{itemize}

\vspace{0.16em}

\compactproject
{基于多智能体 PPO / MAPPO 的智能信号灯协同控制}
{2024.05 -- 2026.05}

\begin{itemize}
  \item 基于 SUMO / TraCI 搭建多路口交通仿真与 PPO / MAPPO 训练框架，并与 FixedTime、MaxPressure、GreenWave 等基线方法进行统一评测。
\end{itemize}

\vspace{0.10em}

\section{获奖情况}

\awarditem
{第一届综合交通运输大模型智能体创新应用大赛，大学生优秀创意奖}
{2025}

\awarditem
{第二十届“九坤杯”北京大学程序设计竞赛，三等奖}
{2022}

\awarditem
{第三十七届全国高中物理竞赛，省级一等奖}
{2020}

\awarditem
{二〇二〇年全国高中数学联赛，省级一等奖}
{2020}

\section{技能}

{\bodyfont
\begin{tabularx}{\linewidth}{>{\heiti}p{3.1cm}X}
编程与工程 & Python，TypeScript / Node.js，C++，Shell，Git，Linux，Docker\\
模型训练与蒸馏 & PyTorch，Trinity-RFT / veRL，vLLM，FSDP，Ray，在线策略蒸馏，Entropy / KL 分析\\
智能体与评测 & Tool Calling / ReAct，Agent Loop，ALFWorld，SWE-bench，Terminal-Bench，自动评测与 JSONL artifacts\\
算法与实验 & Power Sampling，Self-Consistency，PPO / MAPPO，NumPy，Pandas，SUMO / TraCI\\
语言能力 & CET-6，IELTS Academic 6.5\\
\end{tabularx}
}

\end{document}
